\chapter{Infection}

\section*{Patient focus}
Infections can affect almost any part of the body---from the brain to the bones to the skin. Imaging helps your doctor find exactly where an infection is, how severe it is, and whether it needs surgery or just antibiotics. Different scans work better for different body parts, and sometimes more than one scan is needed to get the full picture.


\label{ch:ch57_infection}

\section{Imaging Pattern}
Imaging can localize infection, identify complications, guide sampling or drainage, and monitor response. Common imaging findings of infection include:

\begin{itemize}
  \item \textbf{Edema}: swelling in tissue surrounding the infected area, visible as low attenuation on CT or high signal on T2-weighted MRI
  \item \textbf{Enhancement}: inflamed tissue enhances after contrast due to increased vascularity and capillary leak
  \item \textbf{Restricted diffusion}: pus and abscess contents restrict water movement, producing high signal on DWI---a key differentiator from necrotic tumor
  \item \textbf{Gas}: air within soft tissues or organs suggests gas-forming organisms or fistula
  \item \textbf{Fluid collections}: loculated or free fluid that may represent empyema, abscess, or phlegmon
  \item \textbf{Marrow change}: edema and replacement of normal fatty marrow on MRI in osteomyelitis
  \item \textbf{Lymphadenopathy}: reactive or suppurative nodes adjacent to the infection
\end{itemize}

\section{Infection by Site}

\subsection{Central Nervous System}
Brain abscess appears as a rim-enhancing lesion with central restricted diffusion on MRI. CT may show a hypodense lesion with surrounding edema. Key differentials include necrotic tumor (no restricted diffusion centrally), resolving haematoma, and tumefactive demyelination. MRI with contrast is the investigation of choice.

\subsection{Chest}
Pneumonia presents as consolidation with air bronchograms on chest X-ray and CT. Complications include parapneumonic effusion, empyema (enhancing pleural rind with loculated fluid), lung abscess (cavitating consolidation with air-fluid level), and necrotising pneumonia. Immunocompromised patients may show atypical patterns including ground-glass opacities and micronodules (e.g.\ PCP, fungal infection).

\subsection{Musculoskeletal}
Osteomyelitis shows marrow edema on MRI, cortical destruction, periosteal reaction, and soft tissue extension. Diabetic foot infection requires MRI to assess extent---T1 low signal marrow with corresponding T2 high signal and enhancement. Septic arthritis shows joint effusion, synovial enhancement, periarticular edema, and may progress to cartilage destruction and joint subluxation.

\subsection{Soft Tissue}
Soft tissue abscess appears as a fluid collection with rim enhancement and surrounding edema. Gas within the collection suggests gas-forming organisms. Necrotizing fasciitis shows gas tracking along fascial planes, fluid in the subfascial compartment, and poor definition of fascial boundaries---a surgical emergency.

\subsection{Abdominal}
Pyelonephritis shows renal enlargement, perinephric fat stranding, and nephrogram phase striation on CT. Complications include renal abscess, perinephric abscess, and emphysematous pyelonephritis (gas within renal parenchyma---emergency). Appendicitis, diverticulitis, and cholangitis are common intra-abdominal infections diagnosed on CT or US.

\section{Modality Choice}

\begin{table}[h]
\centering
\caption{Modality selection for infection by site}
\begin{tabular}{|l|l|l|}
\hline
\textbf{Site} & \textbf{First-line} & \textbf{Second-line} \\
\hline
Abdomen & CT with IV contrast & US for biliary/renal \\
Bone & MRI & X-ray (late changes), bone scan \\
Brain & MRI with contrast & CT if MRI unavailable \\
Chest & CXR, then CT & HRCT for immunocompromised \\
Joint & MRI, US & CT for complex anatomy \\
Kidney & US, CT & MRI if contrast allergy \\
Soft tissue & US, MRI & CT for gas detection \\
\hline
\end{tabular}
\end{table}

\section{Nuclear Medicine in Infection}

\textbf{Labelled WBC scan} (Tc-99m or In-111): Useful for osteomyelitis, prosthetic joint infection, and fever of unknown origin. Leukocytes are extracted, labeled, and reinjected; images show focal accumulation at infection sites. Sensitivity is reduced in antibiotic-treated infection.

\textbf{FDG-PET/CT}: High sensitivity for infection in the spine, prosthetic joints, vascular grafts, and fever of unknown origin. FDG accumulates in activated inflammatory cells. Limitation: non-specific---also positive in tumor, post-surgical change, and sterile inflammation.

\section{Distinguishing Infection from Tumor}

\begin{pearlbox}[Infection vs tumor]
Restricted diffusion centrally favours abscess over necrotic tumor. However, high-grade tumors can also restrict. Clinical context, enhancement pattern (smooth thin rim in abscess vs thick irregular rim in tumor), and response to treatment help differentiate. When uncertain, biopsy is indicated.
\end{pearlbox}

\begin{table}[h]
\centering
\caption{Key differentiating features}
\begin{tabular}{|l|l|l|}
\hline
\textbf{Feature} & \textbf{Abscess} & \textbf{Necrotic tumor} \\
\hline
Clinical response & Fever, raised CRP & Weight loss, chronic course \\
DWI center & Restricted (bright) & Usually not restricted \\
Growth rate & Rapid & Variable \\
Rim enhancement & Thin, smooth & Thick, irregular, nodular \\
Surrounding edema & Present & Present \\
\hline
\end{tabular}
\end{table}

\section{Imaging in Immunocompromised Patients}
Immunocompromised patients (HIV, transplant, chemotherapy, biologics) may show:
\begin{itemize}
  \item \textbf{Atypical infections}: PCP (ground-glass, crazy paving), fungal (nodules with halo sign), CMV (diffuse ground-glass)
  \item \textbf{Opportunistic organisms}: Toxoplasmosis, cryptococcosus, PML (JC virus)
  \item \textbf{Blunted inflammatory response}: less enhancement and edema than expected, making diagnosis harder
  \item \textbf{Multiple simultaneous infections}:TB and fungal co-infection, for example
\end{itemize}

HRCT with thin sections is essential. When infection is suspected, recommend biopsy if imaging is non-diagnostic and the patient is stable.

\section{When to Recommend Drainage}
Imaging-guided drainage is indicated when:
\begin{itemize}
  \item A drainable collection $>$3 cm is identified
  \item The patient is failing antibiotic therapy
  \item The collection is causing mass effect or obstruction
  \item A biopsy is needed to distinguish infection from tumor
\end{itemize}
Percutaneous drainage under US or CT guidance is preferred when a safe access route exists. Surgical drainage is reserved for multiloculated collections, inadequately positioned percutaneous drains, or concurrent surgical indication.

\section{Reporting}
State extent, drainable collection, necrosis, bone or joint involvement, vascular complication, and critical anatomy. Correlate with cultures, inflammatory markers, immune status, and treatment response.

\begin{pearlbox}[Reporting checklist]
\begin{enumerate}
  \item Location and extent of infection
  \item Presence and size of drainable collection
  \item Complications (osteonecrosis, empyema, vascular thrombosis, fistula)
  \item Bone/joint involvement and stability
  \item Adjacent critical structures at risk
  \item Comparison with prior imaging when available
\end{enumerate}
\end{pearlbox}

\clearpage
